\lysh{38862}{4}

\begin{alist}
\item Το $4500$ εκφράζει τον συνολικό όγκο του παγετώνα το έτος $2020$, δηλαδή ο συνολικός όγκος ήταν $4500\text{ km}^3$. Το $15$ εκφράζει τον όγκο του πάγου που χάνει κάθε έτος ο παγετώνας, δηλαδή χάνει περίπου $15\text{ km}^3$ πάγου ανά έτος.

\item Για να υπολογίσουμε τον όγκο του παγετώνα το $2040$, αρχικά θα βρούμε το $t$ από τη διαφορά $t = 2040 - 2020 \Leftrightarrow t = 20$. Άρα θα υπολογίσουμε το $V(20)$.
\[V(20) = 4500 - 15 \cdot 20 = 4500 - 300 = 4200\]
Άρα το $2040$ ο όγκος του παγετώνα θα είναι $4200\text{ km}^3$.

\item $V(t) < 500 \Leftrightarrow 4500 - 15t < 500 \Leftrightarrow 4500 - 500 < 15t \Leftrightarrow 4000 < 15t \Leftrightarrow$
\[t > \dfrac{4000}{15} \Leftrightarrow t > 266,6 \simeq 267.\]

Άρα σε $267$ χρόνια μετά το έτος $2020$ ο όγκος του παγετώνα θα πέσει κάτω από $500\text{ km}^3$.
Δηλαδή το έτος $2287$ ($2020+267$) ο όγκος του παγετώνα θα είναι κάτω από $500\text{ km}^3$.

\item
\begin{rlist}
\item Η γραφική παράσταση θα κατασκευαστεί για $0 < t < 300$:
(Σχεδιάζουμε το ευθύγραμμο τμήμα με άκρα τα σημεία $(0, 4500)$ και $(300, 0)$ εξαιρουμένων ίσως των άκρων βάσει του πεδίου ορισμού).
\item Το σημείο $(300,0)$ είναι σημείο της γραφικής παράστασης. Είναι το σημείο τομής της $C$ με τον άξονα του χρόνου $t$.
\textit{Άλλη λύση:}
Επαλήθευση στον τύπο: $V(300) = 4500 - 15 \cdot 300 = 4500 - 4500 = 0$.
\item Το σημείο αυτό δείχνει ότι σε $300$ χρόνια μετά το $2020$, δηλαδή το $2320$, ο παγετώνας θα λιώσει εντελώς.
\end{rlist}
\end{alist}
